% Shared style for the Introduction to Information Geometry lecture slides.
% Compile the lecture .tex files from this directory.
\providecommand{\courseweek}{1}
\usepackage[T1]{fontenc}
\usepackage[utf8]{inputenc}
\usepackage[sc]{mathpazo}
\renewcommand{\sfdefault}{lmss}
\renewcommand{\ttdefault}{lmtt}
\usepackage{amsmath,amssymb,mathtools}
\usepackage{tikz}
\usetikzlibrary{arrows.meta,calc,positioning}
\usepackage{pgfplots}
\pgfplotsset{compat=1.18}
\usepackage{booktabs}

\definecolor{courseink}{HTML}{173749}
\definecolor{coursemuted}{HTML}{526674}
\definecolor{courseteal}{HTML}{0E7287}
\definecolor{coursecopper}{HTML}{B66B43}
\definecolor{courseline}{HTML}{D4E0E3}
\definecolor{coursepale}{HTML}{EFF5F5}

\setbeamercolor{normal text}{fg=courseink,bg=white}
\setbeamercolor{structure}{fg=courseteal}
\setbeamercolor{frametitle}{fg=courseink,bg=white}
\setbeamercolor{title}{fg=courseink}
\setbeamercolor{subtitle}{fg=coursemuted}
\setbeamercolor{block title}{fg=courseink,bg=white}
\setbeamercolor{block body}{fg=courseink,bg=white}
\setbeamercolor{alerted text}{fg=coursecopper}
\setbeamerfont{title}{family=\rmfamily,series=\bfseries,size=\fontsize{32}{37}}
\setbeamerfont{frametitle}{family=\rmfamily,series=\bfseries,size=\fontsize{26}{31}}
\setbeamerfont{normal text}{size=\fontsize{18}{22}}
\setbeamerfont{subtitle}{size=\fontsize{17}{21}}
\setbeamersize{text margin left=0.78cm,text margin right=0.78cm}
\setbeamertemplate{navigation symbols}{}
\setbeamertemplate{itemize items}[circle]
\setbeamertemplate{enumerate items}[default]
\setbeamertemplate{blocks}[default]
\setlength{\parskip}{0.22em}
\setbeamertemplate{footline}{%
  \leavevmode\hbox to \paperwidth{\hspace{.78cm}%
  {\color{coursemuted}\fontsize{8.5}{10}\selectfont Introduction to Information Geometry \hspace{.32cm} / \hspace{.32cm} Week \courseweek}%
  \hfill{\color{coursemuted}\fontsize{8.5}{10}\selectfont \insertframenumber\hspace{.17cm}/\hspace{.17cm}\inserttotalframenumber}\hspace{.78cm}}%
  \vskip.17cm
}
\setbeamertemplate{frametitle}{%
  \vspace{.13cm}
  \insertframetitle\par
  \vspace{.08cm}
}
\setbeamertemplate{title page}{%
  \vspace*{1.00cm}
  {\color{courseteal}\fontsize{10}{13}\selectfont\bfseries
  INTRODUCTION TO INFORMATION GEOMETRY\par}
  \vspace{.40cm}
  {\usebeamerfont{title}\usebeamercolor[fg]{title}\inserttitle\par}
  \vspace{.22cm}
  {\usebeamerfont{subtitle}\usebeamercolor[fg]{subtitle}\insertsubtitle\par}
  \vfill
  {\color{coursemuted}\fontsize{13}{16}\selectfont Christopher Lee\par}
  \vspace{.34cm}
}

\pgfplotsset{
  courseplot/.style={
    width=11.05cm,
    height=5.15cm,
    axis line style={coursemuted!75},
    tick style={coursemuted!65},
    ticklabel style={font=\scriptsize,color=coursemuted},
    label style={font=\normalsize,color=courseink},
    grid=major,
    major grid style={courseline!70,line width=.4pt},
    line width=1.45pt,
    legend style={font=\small,draw=none,fill=none,at={(.98,.03)},anchor=south east},
    clip=true
  }
}

\newcommand{\highlightmath}[1]{{\color{courseteal}#1}}
\newcommand{\thinrule}{\par\vspace{.25cm}{\color{courseline}\rule{\linewidth}{.8pt}}\par}

% A one-dimensional state space.  x=0 and x=1 are placed at 0 and 9 units.
\newcommand{\stateinterval}[1]{%
 \begin{tikzpicture}[x=1cm,y=1cm]
   \draw[courseink,line width=1.25pt] (0,0)--(9,0);
   \foreach \pos/\lab in {0/0,9/1} {
     \fill[courseink] (\pos,0) circle(2.7pt);
     \node[below=9pt,text=coursemuted] at (\pos,0) {\(\lab\)};
   }
   #1
 \end{tikzpicture}%
}

% Paired pages are separate frames in the PDF, with matching titles.
% Supply the whole frame body for each version; the instructor may edit them
% independently without changing overlay logic.
\newcommand{\coursepair}[3]{%
  \begin{frame}{#1}#2\end{frame}%
  \begin{frame}{#1}#3\end{frame}%
}

% Equilateral representation of the categorical simplex, in barycentric
% coordinates with vertices A=(0,0), B=(1,0), C=(1/2,sqrt(3)/2).
\newcommand{\coursesimplex}[1]{%
 \begin{tikzpicture}[scale=4.9]
  \fill[coursepale] (0,0)--(1,0)--(.5,.8660254)--cycle;
  \draw[courseink,line width=.65pt] (0,0)--(1,0)--(.5,.8660254)--cycle;
  \node[below left,font=\small] at (0,0) {A};
  \node[below right,font=\small] at (1,0) {B};
  \node[above,font=\small] at (.5,.8660254) {C};
  #1
 \end{tikzpicture}%
}
